% verification.tex

\section{Verification Strategy}

\begin{frame}{Verification}{Overview}
  \begin{itemize}
    \item 2.1 Energy and emissions accounting
    \item 2.2 Pause / resume logic and hysteresis
    \item 2.3 Progress and runtime consistency
    \item 2.4 Handling of data gaps
    \item 2.5 Unit tests for KPI functions
  \end{itemize}
  \vspace{0.5cm}
  \begin{block}{Purpose}
    Verification checks whether the implementation faithfully follows the
    mathematical model: state updates, decision rules and KPI calculations
    must be correct and internally consistent.
  \end{block}
\end{frame}

% ---------------------------------------------------------------------------
% Detailed example: 2.3 Progress and Runtime Consistency
% ---------------------------------------------------------------------------

\begin{frame}{Verification}{2.3 Progress and Runtime Consistency}
  \textbf{Goal:} Ensure that training progress and wall-clock time evolve
  consistently with the specified baseline training time and pausing behaviour.
  \vspace{0.3cm}

  \textbf{Key state variables}
  \begin{itemize}
    \item $\mathit{training\_progress}$: fraction of total work completed.
    \item $\mathit{elapsed\_wall\_time}$: accumulated wall-clock time.
    \item $\mathit{pause\_count}$: number of pause/resume cycles.
  \end{itemize}
\end{frame}

\begin{frame}{Verification}{2.3 Progress and Runtime Consistency}
  \textbf{Intended update rules}

  For a time step of length $\Delta t$:

  \[
    \Delta p_{\text{step}} =
    \begin{cases}
      \dfrac{\Delta t}{T_{\text{baseline}}} & \text{if not paused} \\
      0 & \text{if paused}
    \end{cases}
  \]

  \[
    \mathit{training\_progress}(t_{k+1}) =
    \mathit{training\_progress}(t_k) + \Delta p_{\text{step}}
  \]

  \[
    \mathit{elapsed\_wall\_time}(t_{k+1}) =
    \mathit{elapsed\_wall\_time}(t_k) + \Delta t
  \]

  Including checkpoint overhead:

  \[
    T_{\text{policy}}^{\star} =
    T_{\text{baseline}} + T_{\text{paused}} +
    \mathit{pause\_count} \cdot \mathit{checkpoint\_overhead\_time}.
  \]
\end{frame}

\begin{frame}{Verification}{2.3 Progress and Runtime Consistency}
  \textbf{Verification procedure (1/2): Baseline case}
  \begin{enumerate}
    \item Disable pausing (or set the threshold so that no pauses occur).
    \item Run the simulation until completion.
    \item Check numerically that:
      \begin{itemize}
        \item $\mathit{training\_progress} \approx 1.0$ at the final time step.
        \item $\mathit{elapsed\_wall\_time} \approx T_{\text{baseline}}$.
      \end{itemize}
    \item Deviations should be within a small numerical tolerance.
  \end{enumerate}
\end{frame}

\begin{frame}{Verification}{2.3 Progress and Runtime Consistency}
  \textbf{Verification procedure (2/2): With pauses}
  \begin{enumerate}
    \item Construct a scenario that triggers a known amount of pausing:
      \begin{itemize}
        \item Measure total paused time $T_{\text{paused}}$.
        \item Record $\mathit{pause\_count}$.
      \end{itemize}
    \item Compute the expected policy runtime
      \[
        T_{\text{policy}}^{\star} =
        T_{\text{baseline}} + T_{\text{paused}} +
        \mathit{pause\_count} \cdot \mathit{checkpoint\_overhead\_time}.
      \]
    \item Compare $T_{\text{policy}}^{\star}$ to the simulated
      $\mathit{elapsed\_wall\_time}$ at completion.
    \item If they match within tolerance, the progress and time accounting are
      consistent with the model.
  \end{enumerate}
\end{frame}